\documentclass[10pt,a4paper]{amsart}
\usepackage{amsmath,amssymb,amsthm,graphicx}
\usepackage{verbatim}
\newenvironment{reference}{\comment}{\endcomment}
%\newenvironment{reference}{}{}
\newenvironment{slogan}{\comment}{\endcomment}
\newenvironment{history}{\comment}{\endcomment}

% For commutative diagrams we use Xy-pic
\usepackage[all]{xy}

% We use 2cell for 2-commutative diagrams.
\xyoption{2cell}
% Theorem environments.
%
\theoremstyle{plain}
\newtheorem{theorem}[subsection]{Theorem}
\newtheorem{proposition}[subsection]{Proposition}
\newtheorem{lemma}[subsection]{Lemma}

\theoremstyle{definition}
\newtheorem{definition}[subsection]{Definition}
\newtheorem{example}[subsection]{Example}
\newtheorem{exercise}[subsection]{Exercise}
\newtheorem{situation}[subsection]{Situation}

\theoremstyle{remark}
\newtheorem{remark}[subsection]{Remark}
\newtheorem{remarks}[subsection]{Remarks}

\numberwithin{equation}{subsection}

% Macros
%
\def\lim{\mathop{\mathrm{lim}}\nolimits}
\def\colim{\mathop{\mathrm{colim}}\nolimits}
\def\Spec{\mathop{\mathrm{Spec}}}
\def\Hom{\mathop{\mathrm{Hom}}\nolimits}
\def\Ext{\mathop{\mathrm{Ext}}\nolimits}
\def\SheafHom{\mathop{\mathcal{H}\!\mathit{om}}\nolimits}
\def\SheafExt{\mathop{\mathcal{E}\!\mathit{xt}}\nolimits}
\def\Sch{\mathit{Sch}}
\def\Mor{\mathop{\mathrm{Mor}}\nolimits}
\def\Ob{\mathop{\mathrm{Ob}}\nolimits}
\def\Sh{\mathop{\mathit{Sh}}\nolimits}
\def\NL{\mathop{N\!L}\nolimits}
\def\CH{\mathop{\mathrm{CH}}\nolimits}
\def\proetale{{pro\text{-}\acute{e}tale}}
\def\etale{{\acute{e}tale}}
\def\QCoh{\mathit{QCoh}}
\def\Ker{\mathop{\mathrm{Ker}}}
\def\Im{\mathop{\mathrm{Im}}}
\def\Coker{\mathop{\mathrm{Coker}}}
\def\Coim{\mathop{\mathrm{Coim}}}

% Boxtimes
%
\DeclareMathSymbol{\boxtimes}{\mathbin}{AMSa}{"02}

%
% Macros for moduli stacks/spaces
%
\def\QCohstack{\mathcal{QC}\!\mathit{oh}}
\def\Cohstack{\mathcal{C}\!\mathit{oh}}
\def\Spacesstack{\mathcal{S}\!\mathit{paces}}
\def\Quotfunctor{\mathrm{Quot}}
\def\Hilbfunctor{\mathrm{Hilb}}
\def\Curvesstack{\mathcal{C}\!\mathit{urves}}
\def\Polarizedstack{\mathcal{P}\!\mathit{olarized}}
\def\Complexesstack{\mathcal{C}\!\mathit{omplexes}}
% \Pic is the operator that assigns to X its picard group, usage \Pic(X)
% \Picardstack_{X/B} denotes the Picard stack of X over B
% \Picardfunctor_{X/B} denotes the Picard functor of X over B
\def\Pic{\mathop{\mathrm{Pic}}\nolimits}
\def\Picardstack{\mathcal{P}\!\mathit{ic}}
\def\Picardfunctor{\mathrm{Pic}}
\def\Deformationcategory{\mathcal{D}\!\mathit{ef}}
\usepackage{cleveref}
\crefname{originallemma}{original lemma}{original lemmas}
\Crefname{originallemma}{Original lemma}{Original lemmas}
\makeatletter
\newlabel{lemma-colimit-sites}{{09YL}{source}{}{sourcepointer-1}{}}
\newlabel{lemma-colimit-sites@cref}{{[originallemma][0][]09YL}{[1][0][]source}{}{}{}}
\makeatother
\crefname{originallemma}{original lemma}{original lemmas}
\Crefname{originallemma}{Original lemma}{Original lemmas}
\makeatletter
\newlabel{categories-lemma-exact-adjoint}{{0039}{source}{}{sourcepointer-2}{}}
\newlabel{categories-lemma-exact-adjoint@cref}{{[originallemma][0][]0039}{[1][0][]source}{}{}{}}
\makeatother
\crefname{originallemma}{original lemma}{original lemmas}
\Crefname{originallemma}{Original lemma}{Original lemmas}
\makeatletter
\newlabel{categories-lemma-directed-commutes}{{002W}{source}{}{sourcepointer-3}{}}
\newlabel{categories-lemma-directed-commutes@cref}{{[originallemma][0][]002W}{[1][0][]source}{}{}{}}
\makeatother
\crefname{originallemma}{original lemma}{original lemmas}
\Crefname{originallemma}{Original lemma}{Original lemmas}
\makeatletter
\newlabel{lemma-limit-sheaf}{{00W2}{source}{}{sourcepointer-4}{}}
\newlabel{lemma-limit-sheaf@cref}{{[originallemma][0][]00W2}{[1][0][]source}{}{}{}}
\makeatother
\title{Computing the pullback of a sheaf along a morphism from a colimit of sites}
\author{The Stacks Project Authors}
\date{stacks-project:sites.tex:local-snapshot-20261012; GNU FDL 1.2 or later}
\begin{document}
\maketitle
\noindent\textit{Editorial note.} All delivered mathematical text is the author's own, copied from the pinned Stacks Project chapter file sites-cohomology.tex; no statement or proof step was rewritten, reordered or removed.\par\medskip
\noindent\textit{Editorial note.} The proof is the author's complete argument; every step the author writes out is inside the counted span and no closing sentence is a placeholder.\par\medskip
\noindent\textit{Editorial note.} Every labelled reference occurring inside the counted spans is carried in assembly.external\_references as a non-load-bearing navigation pointer (same-chapter ones at their definition site, cross-chapter ones at the invocation inside the counted span), so no load-bearing mathematics is replaced by a stub. No bibliography citation occurs inside the counted spans.\par\medskip
\noindent\textit{Editorial note.} The checked chapter list for same-name collisions is the set of stacks chapters staged in this repository: chow, cohomology, divisors, etale-cohomology, flat, more-algebra, more-morphisms, morphisms, spaces-more-morphisms, varieties together with this card's chapter sites-cohomology.\par\medskip
\noindent\textit{Editorial note.} The preamble is the chapter's own retained minimal preamble (sources/stacks-sites-cohomology/retained-preamble.tex), a verbatim excerpt of the native preamble: the Xy-pic setup and the native theorem declarations with the mathematical macros.\par\medskip
\noindent\textit{Allowed standard prerequisites (editorial).}\begin{itemize}
\item The construction of the functors to the colimit site. Same chapter, Lemma lemma-colimit-sites. Supplies the functors u\_i and the continuity checked at the start of the proof.
\item Exactness of adjoints. Categories, Lemma categories-lemma-exact-adjoint. Shows that the pullback, being a left adjoint, is right exact, so that only left exactness has to be proved.
\item Finite limits of sheaves and commutation of filtered colimits with finite limits. Same chapter, Lemma lemma-limit-sheaf, and Categories, Lemma categories-lemma-directed-commutes. Reduce left exactness to the displayed formula, in which filtered colimits commute with finite limits.
\end{itemize}\medskip
\ifdefined\hypertarget\hypertarget{sourcepointer-1}{}\fi
\noindent\textit{Source reference.} Original lemma 09YL: sources/stacks-sites/pinned-original-source.tex, source lines 4079--4079. Navigation pointer invoked at proof line 4079; the counted span contains the pointer only and the cited result is not reproduced inside it, so the pointer is registered as a non-load-bearing navigation stub in the established form for Stacks material. No bibliography citation is involved.\par\medskip
\ifdefined\hypertarget\hypertarget{sourcepointer-2}{}\fi
\noindent\textit{Source reference.} Original lemma 0039: sources/stacks-sites/pinned-original-source.tex, source lines 4096--4096. Navigation pointer invoked at proof line 4096; the counted span contains the pointer only and the cited result is not reproduced inside it, so the pointer is registered as a non-load-bearing navigation stub in the established form for Stacks material. No bibliography citation is involved.\par\medskip
\ifdefined\hypertarget\hypertarget{sourcepointer-3}{}\fi
\noindent\textit{Source reference.} Original lemma 002W: sources/stacks-sites/pinned-original-source.tex, source lines 4054--4054. Navigation pointer invoked at proof line 4054; the counted span contains the pointer only and the cited result is not reproduced inside it, so the pointer is registered as a non-load-bearing navigation stub in the established form for Stacks material. No bibliography citation is involved.\par\medskip
\ifdefined\hypertarget\hypertarget{sourcepointer-4}{}\fi
\noindent\textit{Source reference.} Original lemma 00W2: sources/stacks-sites/pinned-original-source.tex, source lines 4174--4174. Navigation pointer invoked at proof line 4174; the counted span contains the pointer only and the cited result is not reproduced inside it, so the pointer is registered as a non-load-bearing navigation stub in the established form for Stacks material. No bibliography citation is involved.\par\medskip
\begin{situation}
\label{situation-inverse-limit-sites}
Here we are given
\begin{enumerate}
\item a cofiltered index category $\mathcal{I}$,
\item for $i \in \Ob(\mathcal{I})$ a site $\mathcal{C}_i$ such that every
covering in $\mathcal{C}_i$ has a finite index set,
\item for a morphism $a : i \to j$ in $\mathcal{I}$ a morphism of sites
$f_a : \mathcal{C}_i \to \mathcal{C}_j$ given by a continuous functor
$u_a : \mathcal{C}_j \to \mathcal{C}_i$,
\end{enumerate}
such that $f_a \circ f_b = f_c$ whenever $c = a \circ b$ in $\mathcal{I}$.
\end{situation}

\begin{lemma}
\label{lemma-compute-pullback-to-limit}
In Situation \ref{situation-inverse-limit-sites} let
$u_i : \mathcal{C}_i \to \mathcal{C}$ be as constructed in
Lemma \ref{lemma-colimit-sites}. Then $u_i$ defines a morphism
of sites $f_i : \mathcal{C} \to \mathcal{C}_i$. For
$U_i \in \Ob(\mathcal{C}_i)$ and sheaf $\mathcal{F}$ on $\mathcal{C}_i$ we have
\begin{equation}
\label{equation-compute-pullback-to-limit}
f_i^{-1}\mathcal{F}(u_i(U_i)) =
\colim_{a : j \to i} f_a^{-1}\mathcal{F}(u_a(U_i))
\end{equation}
\end{lemma}

\begin{proof}
It is immediate from the arguments in the proof of
Lemma \ref{lemma-colimit-sites} that the functors $u_i$ are continuous.
To finish the proof we have to show that $f_i^{-1} := u_{i, s}$
is an exact functor $\Sh(\mathcal{C}_i) \to \Sh(\mathcal{C})$.
In fact it suffices to show that $f_i^{-1}$ is left exact, because
it is right exact as a left adjoint
(Categories, Lemma \ref{categories-lemma-exact-adjoint}).
We first prove (\ref{equation-compute-pullback-to-limit})
and then we deduce exactness.

\medskip\noindent
For an arbitrary object $V$ of $\mathcal{C}$ we can pick a $a : j \to i$
and an object $V_j \in \Ob(\mathcal{C})$ with $V = u_j(V_j)$. Then we
can set
$$
\mathcal{G}(V) = \colim_{b : k \to j} f_{a \circ b}^{-1}\mathcal{F}(u_b(V_j))
$$
The value $\mathcal{G}(V)$ of the colimit is independent of the choice
of $b : j \to i$ and of the object $V_j$ with $u_j(V_j) = V$; we omit
the verification. Moreover, if $\alpha : V \to V'$ is a morphism of
$\mathcal{C}$, then we can choose $b : j \to i$ and a morphism
$\alpha_j : V_j \to V'_j$ with $u_j(\alpha_j) = \alpha$. This induces
a map $\mathcal{G}(V') \to \mathcal{G}(V)$ by using the restrictions
along the morphisms $u_b(\alpha_j) : u_b(V_j) \to u_b(V'_j)$. A check
shows that $\mathcal{G}$ is a presheaf (omitted).
In fact, $\mathcal{G}$ satisfies the sheaf condition. Namely,
any covering $\mathcal{U} = \{U^t \to U\}$ in $\mathcal{C}$
comes from a finite level. Say $\mathcal{U}_j = \{U^t_j \to U_j\}$
is mapped to $\mathcal{U}$ by $u_j$ for some $a : j \to i$ in $\mathcal{I}$.
Then we have
$$
H^0(\mathcal{U}, \mathcal{G}) =
\colim_{b : k \to j} H^0(u_b(\mathcal{U}_j), f_{b \circ a}^{-1}\mathcal{F}) =
\colim_{b : k \to j} f_{b \circ a}^{-1}\mathcal{F}(u_b(U_j)) =
\mathcal{G}(U)
$$
as desired. The first equality holds because filtered colimits commute
with finite limits
(Categories, Lemma \ref{categories-lemma-directed-commutes}).
By construction $\mathcal{G}(U)$ is given by the right hand side of
(\ref{equation-compute-pullback-to-limit}).
Hence (\ref{equation-compute-pullback-to-limit}) is true if we can
show that $\mathcal{G}$ is equal to $f_i^{-1}\mathcal{F}$.

\medskip\noindent
In this paragraph we check that $\mathcal{G}$ is canonically isomorphic to
$f_i^{-1}\mathcal{F}$. We strongly encourage the reader to skip this paragraph.
To check this we have to show there is a bijection
$\Mor_{\Sh(\mathcal{C})}(\mathcal{G}, \mathcal{H}) =
\Mor_{\Sh(\mathcal{C}_i)}(\mathcal{F}, f_{i, *}\mathcal{H})$
functorial in the sheaf $\mathcal{H}$ on $\mathcal{C}$
where $f_{i, *} = u_i^p$. A map
$\mathcal{G} \to \mathcal{H}$ is the same thing as a compatible
system of maps
$$
\varphi_{a, b, V_j} :
f_{a \circ b}^{-1}\mathcal{F}(u_b(V_j))
\longrightarrow
\mathcal{H}(u_j(V_j))
$$
for all $a : j \to i$, $b : k \to j$ and $V_j \in \Ob(\mathcal{C}_j)$.
The compatibilities force the maps $\varphi_{a, b, V_j}$ to be equal
to $\varphi_{a \circ b, \text{id}, u_b(V_j)}$. Given $a : j \to i$, the
family of maps $\varphi_{a, \text{id}, V_j}$ corresponds to a map
of sheaves $\varphi_a : f_a^{-1}\mathcal{F} \to f_{j, *}\mathcal{H}$.
The compatibilities between the $\varphi_{a, \text{id}, u_a(V_i)}$
and the $\varphi_{\text{id}, \text{id}, V_i}$ implies that $\varphi_a$
is the adjoint of the map $\varphi_{id}$ via
$$
\Mor_{\Sh(\mathcal{C}_j)}(f_a^{-1}\mathcal{F}, f_{j, *}\mathcal{H}) =
\Mor_{\Sh(\mathcal{C}_i)}(\mathcal{F}, f_{a, *}f_{j, *}\mathcal{H}) =
\Mor_{\Sh(\mathcal{C}_i)}(\mathcal{F}, f_{i, *}\mathcal{H})
$$
Thus finally we see that the whole system of maps $\varphi_{a, b, V_j}$
is determined by the map $\varphi_{id} : \mathcal{F} \to f_{i, *}\mathcal{H}$.
Conversely, given such a map $\psi : \mathcal{F} \to f_{i, *}\mathcal{H}$
we can read the argument just given backwards to construct the family
of maps $\varphi_{a, b, V_j}$. This finishes the proof that
$\mathcal{G} = f_i^{-1}\mathcal{F}$.

\medskip\noindent
Assume (\ref{equation-compute-pullback-to-limit}) holds. Then the functor
$\mathcal{F} \mapsto f_i^{-1}\mathcal{F}(U)$ commutes with finite limits
because finite limits of sheaves are computed in the category of presheaves
(Lemma \ref{lemma-limit-sheaf}), the functors $f_a^{-1}$ commutes with finite
limits, and filtered colimits commute with finite limits. To see that
$\mathcal{F} \mapsto f_i^{-1}\mathcal{F}(V)$ commutes with finite limits
for a general object $V$ of $\mathcal{C}$, we can use the same argument using
the formula for $f_i^{-1}\mathcal{F}(V) = \mathcal{G}(V)$ given above.
Thus $f_i^{-1}$ is left exact and the proof of the lemma is complete.
\end{proof}

\end{document}
